\section{Conclusion}
\label{sec:conclusion}

We presented a controlled forward benchmark for neural emulation of solver-generated tsunami-like wavefields. The benchmark pairs procedural bathymetry with diverse initial-displacement sources on shared scenarios, and labels each scenario with up to three numerical references -- hydrostatic reconstruction as the primary shallow-water reference, MUSCL-HR as a second shallow-water reference, and an elevation-only weakly dispersive Boussinesq solver -- so that the same input can be compared across reference choices. On this benchmark we studied FNO and FFNO surrogates against convolutional and operator baselines under a matched protocol, covering accuracy, runtime, distribution shift, cross-resolution transfer, synthetic-source transfer onto fully wet real-bathymetry morphologies, physics diagnostics, and predictive uncertainty.

The core findings are consistent across the diagnostics. Among the single-run direct hydrostatic configurations tested, FFNO is strongest, while the base FNO remains competitive across all three references; the Boussinesq ranking is unit-dependent and reverses in denormalized benchmark-scale RMSE. Single-pass shallow-water models drift upward over the forecast horizon, and seeded windowed variants reduce that error substantially, with Window-FFNO-5 reaching $0.052$ as a first-target-frame-seeded diagnostic scored on the remaining $49$ predicted target frames, not as a like-for-like direct $50$-frame prediction. Inference is several orders of magnitude faster than the CPU NumPy reference implementations used here, which makes large-ensemble and many-scenario studies practical. The strict held-out-family tests are mixed, ranging from mild trench degradation to a rough-source failure, and the 10-crop real-bathymetry suite supports only a modest fully wet morphology-transfer claim under synthetic forcing. The cross-solver diagnostic returns a negative result: the emulator-superiority ratio exceeds one in every pair, so a surrogate reproduces the solver it was trained on far better than it reproduces a different solver, and it does not close the same-scenario, same-saved-index released-label gap between solvers. The deep ensemble gives an uncertainty signal that tracks error usefully (correlation $\approx 0.65$--$0.70$) but is systematically overconfident, with nominal intervals under-covering both in-distribution and under shift.

We emphasize that multi-reference evaluation is what makes these readings trustworthy rather than dependent on a single label choice. Keeping the raw solver-vs-solver gap separate from the surrogate-vs-surrogate error prevents a surrogate's emulation skill from being confused with disagreement between numerical models, and reporting every error against a stated solver and configuration keeps the claims conservative: these are errors against numerical references, not against the real ocean.

Future work is broader. Real-event transfer needs surveyed bathymetry, realistic rupture models, shoreline and wet-dry physics, and validation against observed records. Runtime should be compared with optimized GPU or compiled finite-volume solvers. Full nonlinear Boussinesq references would sharpen cross-reference tests, and uncertainty estimates need recalibration and larger ensembles before any reliability claim. Source reconstruction from sparse observations needs different inputs, metrics, and validation, so we leave it to a separate follow-up study rather than mixing it into the forward-surrogate benchmark.
